% %%%%%%%%%%%%%%%%%%%%%%%%%%%%%%%%%%%%%%%%%%%%%%%%%%%%%%%%%%%%%%%%%%%%%%%%%%%%
%                                Packages import                             %
% %%%%%%%%%%%%%%%%%%%%%%%%%%%%%%%%%%%%%%%%%%%%%%%%%%%%%%%%%%%%%%%%%%%%%%%%%%%%

% Bibliography 
\usepackage{natbib} 

% Typo rules 
\usepackage[french]{babel}        

% Global language 
\selectlanguage{french} 
\usepackage{lipsum}

% Font type class 
\usepackage{mathpazo}  

% --> Changing section font
\usepackage{sectsty}                        
\allsectionsfont{\sffamily}

% Source input encoding 
\usepackage[utf8]{inputenc}    

% Source output encoding
\usepackage[T1]{fontenc}    

% For typographic adjustments
\usepackage{microtype}    

% Caption settings
\usepackage[margin=17pt,font=small,labelfont={bf,sf}]{caption}     

% Margins and page layout             
\usepackage[verbose,tmargin=3cm,bmargin=3cm,
            lmargin=2cm,rmargin=2cm,headheight=3cm,
            headsep=0.5cm,footskip=0.8cm]{geometry}   
            
% Interline settings            
\usepackage{setspace}                       
\setstretch{1.1}

% Comment package
\usepackage{verbatim} 

% Table packages
\usepackage{multirow}                       
\usepackage{multicol}
\usepackage{arydshln}
\usepackage{booktabs}

% Figure placement
\usepackage{float}                          
\usepackage{graphicx}
\usepackage{adjustbox}
\usepackage{lscape}

% Hyperlinks  
\usepackage{url}                            
\usepackage[breaklinks=true,unicode=true,colorlinks=true,
            linkcolor=tablemat,citecolor=citation,
            urlcolor=Blue,backref=page]{hyperref}

% Math facilities            
\usepackage{amsmath}                        

% Cross-references settings
\usepackage{cleveref}                       
\crefname{figure}{Fig.}{Fig.}
\crefname{table}{Tab.}{Tab.}
\crefname{section}{Sec.}{Sec.}
\crefname{equation}{Eq.}{Eq.}
\crefname{lstlisting}{Ex.}{Ex.}

\newcommand{\caw}{\href{https://gitlab.com/cawaqs/cawaqs}{\texttt{CaWaQS}\ }} 

% To fill pages from the top
\raggedbottom

% Algorithm pkg
\usepackage[]{algorithm2e}
\usepackage{algpseudocode}
\renewcommand{\algorithmcfname}{Algorithm}  

% Symbols list  
\usepackage{nomencl}                         

% Itemize settings
\usepackage{enumitem}                       
\setlist[itemize]{label=$\bullet$, nosep, 
        after=\vspace{0.2cm}}               
\setlist{noitemsep}

\usepackage{etoolbox}
\makenomenclature

% Color options
\usepackage{color}                          
\usepackage[dvipsnames,svgnames,table]{xcolor}
\definecolor{tablemat}{rgb}{0.18,0.32,0.32}
\definecolor{citation}{rgb}{0,0,0.45}
\definecolor{darkWhite}{rgb}{0.98,0.98,0.98}
\definecolor{darkGreen}{rgb}{0,0.5,0}

% Lst for COMM and input file examples
\usepackage{listings}                       

% Custom boxe fFor warning, info, etc. 
\usepackage{Setup/Messages_pkg/alertmessage}    % ** WARNING ** Don't switch with xcolor --> conflict otherwise !

% Alertmessage pkg command memo
% \alertinfo{Test.}
% \alertsuccess{Test.}
% \alertwarning{Test.}
% \alerterror{Test.}

% %%%%%%%%%%%%%%%%%%%%%%%%%%%%%%%%%%%%%%%%%%%%%%%%%%%%%%%%%%%%%%%%%%%%%%%%%%%%
%                         Document layout settings                           %
% %%%%%%%%%%%%%%%%%%%%%%%%%%%%%%%%%%%%%%%%%%%%%%%%%%%%%%%%%%%%%%%%%%%%%%%%%%%%

% General headers and page bottom settings 
\usepackage{fancyhdr}                       
\pagestyle{fancy}
\fancyhead[LE]{\thepage}
\fancyhead[RE]{\textbf{\leftmark}}
\fancyhead[RO]{\thepage}
\fancyhead[LO]{\textbf{\rightmark}}
\fancyfoot[L,C,R]{}

% Chapter header settings 
\usepackage[Bjornstrup]{fncychap}  
\ChNumVar{\Large{\textbf{\color{white}{CHAPTER}}} \fontsize{76}{80}\usefont{OT1}{pzc}{m}{n}\selectfont}
\ChTitleVar{\raggedright\LARGE\bfseries} 
\addto\captionsfrench{\def\tablename{Table}} 

% Table of contents settings
\usepackage{tocloft}     

\renewcommand{\cftchapfont}{\scshape}
\renewcommand{\cftsecfont}{\bfseries}
\renewcommand{\cftfigfont}{Figure }
\renewcommand{\cfttabfont}{Table }
\setcounter{secnumdepth}{3}
\setcounter{tocdepth}{4}

% Minitoc settings 
\usepackage[french]{minitoc}
\mtcsettitle{minitoc}{Content}
\mtcsetrules{*}{on}
\mtcsetoffset{minitoc}{0pt}
\mtcsetdepth{minitoc}{4}
\mtcsetpagenumbers{minitoc}{on}
\mtcsetformat{minitoc}{tocrightmargin}{1mm}
\mtcsetformat{minitoc}{dotinterval}{0.11mm}
\makeatletter
\renewcommand{\@pnumwidth}{4mm}
\makeatother
\providecommand{\tabularnewline}{\\}
\setlength{\mtcindent}{0cm} % No indent in miniTOC

\renewcommand{\cftfigfont}{\raggedright} % Align figure titles to the left
\renewcommand{\cftfignumwidth}{3em}      % Adjust space for figure number
\renewcommand{\cfttabfont}{\raggedright} 
\renewcommand{\cfttabnumwidth}{3em}

% More chapter header settings
% Just changing color here to match the TOC colors. 

% Chapter title color
\newcommand{\colortitlechap}{\color[rgb]{1,1,1}} 

% Chapter number color
\newcommand{\colornumberchap}{\color[rgb]{0.7,0.8,0.2}} 

% Background box (= tablemat color) 
\newcommand{\colorbackchap}{\colorbox[rgb]{0.18,0.32,0.32}} 

% Chapter headers
\makeatletter
\renewcommand{\DOCH}{%
    \settowidth{\py}{\CNoV\thechapter}
    \addtolength{\py}{-10pt}% 
    \fboxsep=0pt%
    \colorbackchap{\rule{0pt}{40pt}\parbox[b]{\textwidth}{\hfill}}%
    \kern-\py\raise20pt%
    \hbox{\colornumberchap\CNoV\thechapter}\\%
}

\renewcommand{\DOTI}[1]{%
    \nointerlineskip\raggedright%
    \fboxsep=\myhi%
    \vskip-1ex%
    \colorbackchap{\parbox[t]{\mylen}{\CTV\FmTi{\colortitlechap#1}}}\par\nobreak%
    \vskip 40\p@%
}

\renewcommand{\DOTIS}[1]{%
    \fboxsep=0pt%
    \colorbackchap{\rule{0pt}{40pt}\parbox[b]{\textwidth}{\hfill}}\\%
    \nointerlineskip\raggedright%
    \fboxsep=\myhi%
    \colorbackchap{\parbox[t]{\mylen}{\CTV\FmTi{\colortitlechap#1}}}\par\nobreak%
    \vskip 40\p@%
}

\makeatother

% Settings for COMM/code/command excerpts
\lstset{aboveskip=6mm,belowskip=0mm,backgroundcolor=\color{darkWhite}, basicstyle=\ttfamily\footnotesize,breakatwhitespace=false,
 breaklines=true,captionpos=b,commentstyle=\color{blue},deletekeywords={...},escapeinside={\%*}{*)},extendedchars=true, framexleftmargin=0pt,
  framextopmargin=2pt,framexbottommargin=2pt,frame=tb, keepspaces=true, keywordstyle=\color{blue}, language=C,morekeywords={*,...},
  numbers=left,numbersep=3pt,numberstyle=\tiny\color{black},rulecolor=\color{black},showspaces=false,showstringspaces=false,showtabs=false,stepnumber=1, 
  stringstyle=\color{gray}, tabsize=3, title=\lstname, morecomment=[l][\color{blue}]{!}}

% Listings caption settings
\renewcommand{\lstlistingname}{Example} 
\renewcommand{\lstlistlistingname}{List of file, code and command examples}

% insert python code in latex
\usepackage{minted}

\setminted{
  fontsize=\normal,
  breaklines=true,
  linenos=false,
  frame=lines,
  framesep=2mm
}

% For \foreach loops in author/affiliation macros
\usepackage{pgffor}

% --- Compteurs ---
\newcounter{nid}          % Compteur pour les auteurs
\newcounter{nrecord}      % Compteur pour les affiliations

% --- Macro pour ajouter un auteur (format stocké : Prénoms Nom) ---
\newcommand{\addauthor}[2]{%
    \stepcounter{nid}%
    \expandafter\xdef\csname authname\arabic{nid}\endcsname{#2 #1}%
}

% --- Macro pour ajouter une affiliation ---
\newcommand{\addaffiliation}[1]{%
    \stepcounter{nrecord}%
    \expandafter\xdef\csname affilitext\arabic{nrecord}\endcsname{#1}%
}

% --- Macro pour associer un auteur à une ou plusieurs affiliations ---
% Usage: \setaffiliations{1,2,3} after \addauthor
\newcommand{\setaffiliations}[1]{%
    \expandafter\xdef\csname authaffils\arabic{nid}\endcsname{#1}%
}

% --- Macro pour lister les auteurs avec leurs numéros d'affiliation en exposant ---
\newcommand{\listAuthors}{%
    \foreach \i in {1,...,\value{nid}} {%
        \ifnum\i>1, \fi%
        \csname authname\i\endcsname%
        \textsuperscript{\csname authaffils\i\endcsname}%
    }%
}

% --- Macro pour lister uniquement les noms des auteurs (sans numéros) ---
\newcommand{\listAuthorNames}{%
    \foreach \i in {1,...,\value{nid}} {%
        \ifnum\i>1, \fi%
        \csname authname\i\endcsname%
    }%
}

% --- Macro pour lister toutes les affiliations numérotées ---
\newcommand{\listAffiliations}{%
    \foreach \i in {1,...,\value{nrecord}} {%
        \textsuperscript{\i}~\csname affilitext\i\endcsname\\[0.1cm]%
    }%
}

% --- Macro pour lister les affiliations d'un auteur spécifique ---
\newcommand{\listAuthorAffiliations}[1]{%
    \edef\affilids{\csname authaffils#1\endcsname}%
    \foreach \id in \affilids {%
        \textsuperscript{\id}~\csname affilitext\id\endcsname\\%
    }%
}

% --- Macro pour lister tous les auteurs avec leurs affiliations ---
\newcommand{\listAllAuthorsWithAffiliations}{%
    \foreach \i in {1,...,\value{nid}} {%
        \textbf{\csname authname\i\endcsname}%
        \textsuperscript{\csname authaffils\i\endcsname}%
        \\[0.2cm]%
    }%
}
